\documentclass[a4paper,12pt]{article}
\usepackage[utf8]{inputenc}
\usepackage[T1]{fontenc}
\usepackage[french]{babel}
\usepackage{amsmath,amsfonts,amssymb}
\usepackage{graphicx}
\usepackage{hyperref}
\usepackage{listings}
\usepackage{xcolor}
\usepackage{geometry}
\geometry{hmargin=2.5cm, vmargin=2.5cm}

% Configuration de l'affichage du code
\definecolor{codegreen}{rgb}{0,0.6,0}
\definecolor{codegray}{rgb}{0.5,0.5,0.5}
\definecolor{codepurple}{rgb}{0.58,0,0.82}
\definecolor{backcolour}{rgb}{0.95,0.95,0.92}

\lstdefinestyle{mystyle}{
    backgroundcolor=\color{backcolour},
    commentstyle=\color{codegreen},
    keywordstyle=\color{magenta},
    numberstyle=\tiny\color{codegray},
    stringstyle=\color{codepurple},
    basicstyle=\ttfamily\small,
    breakatwhitespace=false,
    breaklines=true,
    captionpos=b,
    keepspaces=true,
    numbers=left,
    numbersep=5pt,
    showspaces=false,
    showstringspaces=false,
    showtabs=false,
    tabsize=4,
    language=Python,
    literate={é}{{\'e}}1 {è}{{\`e}}1 {ê}{{\^e}}1 {à}{{\`a}}1 {â}{{\^a}}1 {ç}{{\c{c}}}1,
}

\lstset{style=mystyle}

\title{\textbf{Cours d'Informatique pour Débutants} \\
\large Programmation Orientée Objet en Python \& Conteneurisation avec Docker}
\author{Votre Enseignant}
\date{\today}

\begin{document}

\maketitle
\tableofcontents
\newpage

% ======================================================================
% PREMIER COURS : PROGRAMMATION ORIENTEE OBJET EN PYTHON
% ======================================================================
\section{Programmation Orientée Objet en Python}

\subsection{Introduction à la Programmation Orientée Objet (POO)}
La programmation orientée objet est un paradigme de programmation qui organise le code autour de la notion d'\textbf{objets}. Un objet est une entité qui possède des \textbf{attributs} (ses caractéristiques) et des \textbf{méthodes} (ses actions). Cette approche permet de modéliser le monde réel de manière plus intuitive.

\subsection{Différence entre Programmation Procédurale et Orientée Objet}
\begin{itemize}
    \item \textbf{Procédurale} : On se concentre sur les fonctions et la logique. Les données et les fonctions sont séparées. (Exemple : C, Pascal).
    \item \textbf{Orientée Objet} : On regroupe les données (attributs) et les fonctions (méthodes) qui agissent sur ces données dans des \textbf{classes}. (Exemple : Python, Java, C++).
\end{itemize}

\subsection{Définition d'une classe en Python}
Une classe est un modèle (un moule) qui permet de créer des objets.
\begin{lstlisting}
class Chien:
    # Le corps de la classe
    pass
\end{lstlisting}

\subsection{Création d'objets}
Un objet est une \textbf{instance} d'une classe.
\begin{lstlisting}
class Chien:
    pass

# Creation d'objets (instances)
mon_chien = Chien()
ton_chien = Chien()

print(mon_chien) # Affiche : <__main__.Chien object at 0x...>
\end{lstlisting}

\subsection{Attributs et Méthodes}
\begin{itemize}
    \item \textbf{Attributs} : Variables propres à l'objet.
    \item \textbf{Méthodes} : Fonctions propres à l'objet.
\end{itemize}
\begin{lstlisting}
class Chien:
    # Methode
    def aboyer(self):
        print("Woof !")

# Creation
mon_chien = Chien()
mon_chien.nom = "Rex"  # Ajout dynamique d'un attribut (peu recommande)
mon_chien.aboyer()     # Appel de la methode
\end{lstlisting}

\subsection{Le constructeur \texttt{\_\_init\_\_}}
Le constructeur est une méthode spéciale appelée automatiquement à la création de l'objet. Il sert à initialiser les attributs.
\begin{lstlisting}
class Chien:
    def __init__(self, nom, age):
        self.nom = nom   # Attribut d'instance
        self.age = age   # Attribut d'instance

    def aboyer(self):
        return f"{self.nom} dit Woof !"

# Utilisation
mon_chien = Chien("Medor", 5)
print(mon_chien.aboyer())
\end{lstlisting}

\subsection{Encapsulation}
L'encapsulation est le principe qui consiste à restreindre l'accès direct à certaines données. En Python, on utilise une convention : un attribut précédé d'un underscore \texttt{\_} est considéré comme "protégé".
\begin{lstlisting}
class CompteBancaire:
    def __init__(self, titulaire, solde):
        self.titulaire = titulaire
        self._solde = solde  # Attribut protege

    def deposer(self, montant):
        if montant > 0:
            self._solde += montant

    def afficher_solde(self):
        return f"Solde : {self._solde} euros"
\end{lstlisting}

\subsection{Méthodes \texttt{get} et \texttt{set}}
On utilise des méthodes pour accéder (getter) et modifier (setter) les attributs protégés de manière contrôlée.
\begin{lstlisting}
class Personne:
    def __init__(self, nom):
        self._nom = nom

    # Getter
    def get_nom(self):
        return self._nom

    # Setter
    def set_nom(self, nouveau_nom):
        if nouveau_nom:
            self._nom = nouveau_nom
        else:
            print("Erreur : nom invalide")

p = Personne("Alice")
print(p.get_nom())
p.set_nom("Bob")
\end{lstlisting}

\subsection{Héritage}
L'héritage permet de créer une nouvelle classe (classe fille) à partir d'une classe existante (classe mère). La fille hérite des attributs et méthodes de la mère.
\begin{lstlisting}
class Animal: # Classe mere
    def __init__(self, nom):
        self.nom = nom
    def parler(self):
        return "..."

class Chat(Animal): # Classe fille
    def parler(self):
        return f"{self.nom} dit Miaou !"

class Chien(Animal): # Classe fille
    def parler(self):
        return f"{self.nom} dit Woof !"

# Utilisation
animaux = [Chat("Felix"), Chien("Rex")]
for a in animaux:
    print(a.parler())
\end{lstlisting}

\subsection{Polymorphisme}
Le polymorphisme (du grec "plusieurs formes") permet à différentes classes d'utiliser une même méthode de manière spécifique. Dans l'exemple ci-dessus, la méthode \texttt{parler()} est polymorphique.

\subsection{Exemple complet de programme Python}
\begin{lstlisting}
# Definition d'une classe Vehicule
class Vehicule:
    def __init__(self, marque, modele, annee):
        self.marque = marque
        self.modele = modele
        self.annee = annee
        self._kilometrage = 0

    def rouler(self, km):
        if km > 0:
            self._kilometrage += km
            print(f"Le vehicule a roule {km} km.")
        else:
            print("Distance invalide.")

    def get_kilometrage(self):
        return self._kilometrage

    def __str__(self):
        return f"{self.marque} {self.modele} ({self.annee})"

# Heritage : Voiture est une sorte de Vehicule
class Voiture(Vehicule):
    def __init__(self, marque, modele, annee, nb_portes):
        super().__init__(marque, modele, annee) # Appel du constructeur parent
        self.nb_portes = nb_portes

    def klaxonner(self):
        print("Tut tut !")

    def __str__(self):
        return super().__str__() + f" - {self.nb_portes} portes"

# Programme principal
if __name__ == "__main__":
    ma_voiture = Voiture("Toyota", "Yaris", 2020, 5)
    print(ma_voiture)
    ma_voiture.rouler(150)
    print(f"Kilometrage: {ma_voiture.get_kilometrage()} km")
    ma_voiture.klaxonner()
\end{lstlisting}

\subsection{Exercices Corrigés}
\subsubsection{Exercice 1 : Livre}
Créez une classe \texttt{Livre} avec les attributs \texttt{titre}, \texttt{auteur} et \texttt{annee}. Ajoutez une méthode \texttt{description()} qui retourne une chaîne formatée.
\\[0.5cm]
\textbf{Correction :}
\begin{lstlisting}
class Livre:
    def __init__(self, titre, auteur, annee):
        self.titre = titre
        self.auteur = auteur
        self.annee = annee
    def description(self):
        return f"'{self.titre}' par {self.auteur} ({self.annee})"

livre = Livre("1984", "George Orwell", 1949)
print(livre.description())
\end{lstlisting}

\subsubsection{Exercice 2 : Compteur}
Créez une classe \texttt{Compteur} avec un attribut privé \texttt{\_valeur}. Implémentez les méthodes \texttt{incrementer()} et \texttt{get\_valeur()}. Empêchez la valeur de devenir négative.
\\[0.5cm]
\textbf{Correction :}
\begin{lstlisting}
class Compteur:
    def __init__(self):
        self._valeur = 0
    def incrementer(self, pas=1):
        if pas > 0:
            self._valeur += pas
    def get_valeur(self):
        return self._valeur

c = Compteur()
c.incrementer(5)
print(c.get_valeur())
\end{lstlisting}

\newpage

% ======================================================================
% DEUXIEME COURS : CONTENEURISATION AVEC DOCKER
% ======================================================================
\section{Conteneurisation avec Docker}

\subsection{Introduction à la conteneurisation}
La conteneurisation est une méthode de virtualisation légère qui permet d'empaqueter une application et ses dépendances dans un environnement isolé appelé \textbf{conteneur}. Cela garantit que l'application fonctionne de manière identique quel que soit l'environnement (développement, test, production).

\subsection{Différence entre Machines Virtuelles et Conteneurs}
\begin{itemize}
    \item \textbf{Machine Virtuelle (VM)} : Inclut le système d'exploitation invité complet, ce qui la rend lourde (plusieurs Go) et lente à démarrer. L'hyperviseur virtualise le matériel.
    \item \textbf{Conteneur} : Partage le noyau du système d'exploitation hôte. Il est léger (quelques Mo), démarre en quelques secondes et consomme moins de ressources.
\end{itemize}

\subsection{Présentation de Docker}
Docker est la plateforme de conteneurisation la plus populaire. Elle permet de créer, déployer et exécuter des applications dans des conteneurs.

\subsection{Installation de Docker}
\begin{itemize}
    \item \textbf{Windows/Mac} : Télécharger Docker Desktop sur \url{https://www.docker.com/products/docker-desktop/}
    \item \textbf{Linux (Ubuntu/Debian)} :
    \begin{lstlisting}[language=bash]
sudo apt update
sudo apt install docker.io
sudo systemctl start docker
sudo systemctl enable docker
    \end{lstlisting}
\end{itemize}

\subsection{Architecture de Docker}
\begin{itemize}
    \item \textbf{Docker Engine} : Le moteur qui construit et exécute les conteneurs.
    \item \textbf{Images} : Modèles en lecture seule (comme une classe en POO) utilisés pour créer des conteneurs.
    \item \textbf{Conteneurs} : Instances exécutables d'une image (comme un objet).
    \item \textbf{Registry} : Dépôt d'images (ex: Docker Hub).
\end{itemize}

\subsection{Explication détaillée des commandes Docker}

\subsubsection{\texttt{docker pull}}
Télécharge une image depuis un registry sans l'exécuter.
\begin{lstlisting}[language=bash]
docker pull hello-world
docker pull python:3.9-slim  # Image Python officielle
\end{lstlisting}

\subsubsection{\texttt{docker images}}
Liste les images téléchargées sur votre machine.
\begin{lstlisting}[language=bash]
docker images
\end{lstlisting}

\subsubsection{\texttt{docker run}}
Crée et démarre un conteneur à partir d'une image.
\begin{itemize}
    \item \texttt{-d} : mode détaché (en arrière-plan)
    \item \texttt{-it} : mode interactif avec terminal
    \item \texttt{--name} : donne un nom au conteneur
    \item \texttt{-p 8080:80} : mappe le port 8080 de l'hôte au port 80 du conteneur
\end{itemize}
\begin{lstlisting}[language=bash]
docker run hello-world
docker run -d -p 8080:80 --name mon_nginx nginx
docker run -it python:3.9-slim bash
\end{lstlisting}

\subsubsection{\texttt{docker ps}}
Liste les conteneurs en cours d'exécution. \texttt{docker ps -a} liste tous les conteneurs (même arrêtés).
\begin{lstlisting}[language=bash]
docker ps
docker ps -a
\end{lstlisting}

\subsubsection{\texttt{docker stop}}
Arrête un ou plusieurs conteneurs en cours d'exécution.
\begin{lstlisting}[language=bash]
docker stop mon_nginx
docker stop id_conteneur
\end{lstlisting}

\subsubsection{\texttt{docker start}}
Démarre un conteneur arrêté.
\begin{lstlisting}[language=bash]
docker start mon_nginx
\end{lstlisting}

\subsubsection{\texttt{docker rm}}
Supprime un ou plusieurs conteneurs (ils doivent être arrêtés).
\begin{lstlisting}[language=bash]
docker rm mon_nginx
docker rm id_conteneur
\end{lstlisting}

\subsubsection{\texttt{docker exec}}
Exécute une commande dans un conteneur en cours d'exécution.
\begin{lstlisting}[language=bash]
docker exec -it mon_nginx bash  # Ouvre un shell bash dans le conteneur
docker exec mon_nginx ls -l     # Liste les fichiers sans entrer dans le shell
\end{lstlisting}

\subsubsection{\texttt{docker build}}
Construit une image à partir d'un fichier \texttt{Dockerfile}.
\begin{lstlisting}[language=bash]
docker build -t mon_app:latest .
# -t donne un nom et un tag a l'image
# . indique le chemin du contexte de build (souvent le dossier courant)
\end{lstlisting}

\subsection{Création d'un conteneur Docker étape par étape}
\begin{enumerate}
    \item Télécharger une image : \texttt{docker pull nginx}
    \item Créer et démarrer un conteneur : \texttt{docker run -d --name site\_web -p 8080:80 nginx}
    \item Vérifier qu'il tourne : \texttt{docker ps}
    \item Ouvrir un navigateur à l'adresse \texttt{http://localhost:8080}
    \item Arrêter le conteneur : \texttt{docker stop site\_web}
    \item Le supprimer : \texttt{docker rm site\_web}
\end{enumerate}

\subsection{Création d'un Dockerfile}
Un \texttt{Dockerfile} est un fichier texte contenant les instructions pour construire une image Docker.

\subsubsection{Instructions détaillées}
\begin{itemize}
    \item \textbf{FROM} : Spécifie l'image de base. \texttt{FROM python:3.9-slim}
    \item \textbf{WORKDIR} : Définit le répertoire de travail dans le conteneur. \texttt{WORKDIR /app}
    \item \textbf{COPY} : Copie des fichiers de l'hôte vers le conteneur. \texttt{COPY requirements.txt .}
    \item \textbf{RUN} : Exécute des commandes lors de la construction de l'image (ex: installation de dépendances). \texttt{RUN pip install -r requirements.txt}
    \item \textbf{CMD} : Commande exécutée par défaut quand le conteneur démarre. Il ne peut y en avoir qu'une. \texttt{CMD ["python", "app.py"]}
    \item \textbf{EXPOSE} : Indique le port sur lequel le conteneur écoute (informationnelle). \texttt{EXPOSE 5000}
\end{itemize}

\subsubsection{Exemple complet : Dockerfile pour une application Python}
Imaginons une application Python simple avec Flask.
\\[0.5cm]
\textbf{Structure du projet :}
\begin{verbatim}
mon_app/
    app.py
    requirements.txt
    Dockerfile
\end{verbatim}
\textbf{Contenu de \texttt{app.py}:}
\begin{lstlisting}
from flask import Flask
app = Flask(__name__)

@app.route('/')
def hello():
    return "Hello, Docker !"

if __name__ == '__main__':
    app.run(host='0.0.0.0', port=5000)
\end{lstlisting}
\textbf{Contenu de \texttt{requirements.txt}:}
\begin{lstlisting}
Flask==2.2.3
\end{lstlisting}
\textbf{Contenu du \texttt{Dockerfile}:}
\begin{lstlisting}[language=bash]
# Utiliser une image Python officielle comme base
FROM python:3.9-slim

# Definir le repertoire de travail
WORKDIR /app

# Copier le fichier des dependances
COPY requirements.txt .

# Installer les dependances Python
RUN pip install --no-cache-dir -r requirements.txt

# Copier tout le code de l'application
COPY . .

# Indiquer le port utilise par l'application
EXPOSE 5000

# Commande pour lancer l'application
CMD ["python", "app.py"]
\end{lstlisting}
\textbf{Construction et exécution :}
\begin{lstlisting}[language=bash]
docker build -t ma_flask_app .
docker run -d -p 5000:5000 --name mon_site ma_flask_app
# Aller sur http://localhost:5000
\end{lstlisting}

\subsection{Exercices Pratiques avec Correction}
\subsubsection{Exercice 1 : Conteneuriser un script simple}
Créez un script Python \texttt{hello.py} qui affiche "Bonjour le monde !". Écrivez un \texttt{Dockerfile} basé sur \texttt{python:alpine} qui copie le script et l'exécute avec \texttt{CMD}.
\\[0.5cm]
\textbf{Correction :}
\begin{lstlisting}[language=bash]
# hello.py
print("Bonjour le monde !")

# Dockerfile
FROM python:alpine
WORKDIR /app
COPY hello.py .
CMD ["python", "hello.py"]

# Construction et execution
docker build -t hello .
docker run hello
\end{lstlisting}

\subsubsection{Exercice 2 : Gestion de conteneurs}
1. Lancez un conteneur \texttt{nginx} nommé \texttt{web1} sur le port 8081. \\
2. Vérifiez qu'il tourne. \\
3. Ouvrez un shell à l'intérieur et modifiez la page d'accueil (fichier \texttt{/usr/share/nginx/html/index.html}). \\
4. Arrêtez et supprimez le conteneur.
\\[0.5cm]
\textbf{Correction :}
\begin{lstlisting}[language=bash]
# 1.
docker run -d --name web1 -p 8081:80 nginx
# 2.
docker ps
# 3.
docker exec -it web1 bash
echo "<h1>Ma page modifiee</h1>" > /usr/share/nginx/html/index.html
exit
# 4.
docker stop web1
docker rm web1
\end{lstlisting}

\end{document}